\section{Introduction}
\label{sec:introduction}

\subsection{Background and Motivation}
\label{sec:background}

Modern offline Handwritten Text Recognition (HTR) has matured on the widely used benchmarks. A compact convolutional--recurrent CTC network with careful preprocessing~\cite{retsinas2024htrbestpractices} reaches roughly $4.6\%$ character error rate (CER) on IAM~\cite{marti2002iam}, and Transformer-based encoders such as HTR-VT~\cite{li2024htrvt} and pretrained encoder--decoder systems such as TrOCR~\cite{li2023trocr} reach comparable or lower error at higher parameter cost. What remains unresolved is what happens inside these recognizers. Downstream document-analysis tasks such as writer identification, style analysis, and palaeographic dating have traditionally relied on hand-designed features~\cite{raven2025vlac}, and a natural alternative is to reuse the internal attention of a Transformer HTR system directly. This alternative requires the attention maps to be clean, that is, peaks that shift with the recognized character, low background diffusion, and no spatially arbitrary artifacts.

The obstacle to reusing ViT attention was documented by Darcet~\etal~\cite{darcet2024registers}. In supervised and self-supervised Vision Transformers, a small number of patch tokens develop anomalously high norms and absorb a disproportionate share of the attention mass. These \emph{attention sinks} produce visible artifacts in attention maps and hurt any downstream task that reads those maps. Their proposed fix is minimally invasive: prepend a small number of learnable \emph{register tokens} to the patch sequence, and the sinks migrate onto the registers, leaving patch attention interpretable. On classification ViTs the fix is clean; it improves downstream detection and segmentation while leaving top-1 accuracy unchanged. The mechanism has not been evaluated for CTC-based sequence-labelling HTR, where the queries that read the attention are the patch tokens themselves and no CLS-style aggregate query exists.

This project sits at that gap. The concrete deliverable is a modular HTR pipeline that trains a CNN-stem Vision Transformer with configurable register tokens (ViT-RGTS), keeps the CNN--RNN--CTC baseline of Retsinas~\etal\ as the reference point, and adds pretrained lanes (TorchVision ViT-B/16, TrOCR-base) for a fair architectural comparison. The concrete question is whether register tokens, evaluated at $R \in \{0, 2, 4, 8, 16\}$ on IAM and READ2016, change (a) the character/word error rate and (b) the quality of the attention maps produced by the encoder, judged both qualitatively in the style of Fig.~5 of \emph{Beyond Memorization}~\cite{gurav2025beyondmem} and quantitatively through attention-concentration metrics.

\subsection{Problem Statement and Scope}
\label{sec:problem}

The project builds a Vision Transformer HTR system whose attention maps are fine-grained enough to support downstream writer-analysis tasks. Register tokens are the specific mechanism under test, motivated by the classification-ViT results of Darcet~\etal~\cite{darcet2024registers}. Success is measured jointly on character/word error rate and on attention quality; the writer-identification pipeline itself is out of scope for a 10~ECTS project budget and is sketched as future work.

The report contributes on five fronts. First, it shows that register tokens are CER-neutral within seed noise on both IAM and READ2016, replicating on a CTC sequence recognizer what Darcet~\etal\ reported for classification ViTs. Second, it shows on a compact set of attention-quality metrics that register tokens reduce character-attention entropy, improve the spatial ordering of the peaks, and recruit more heads into the left-to-right reading pattern, while a residual sink persists on patch tokens. Third, it reproduces the per-character attention visualization of Gurav~\etal\ end-to-end on the register-added ViT. Fourth, it identifies the CNN stem, and the strict left-to-right token order it enforces, as a structural requirement for CTC ViTs: two independent architectural configurations that violate this order collapse to nearly $73\%$ CER. Fifth, it releases the complete training and evaluation pipeline together with the trained models and result artifacts, so that the writer-identification step suggested by the problem statement can be built directly on top.

Section~\ref{sec:related_work} surveys the HTR, register-token, and attention-interpretability literature and positions the present work within it. Section~\ref{sec:methodology} describes the pipeline, datasets, and models. Section~\ref{sec:experiments} lays out the evaluation protocol. Section~\ref{sec:results} presents the recognition results, reproduces the \emph{Beyond Memorization} visualization on the trained model, and critically assesses whether the Darcet~\etal\ claim holds in the CTC HTR setting. Section~\ref{sec:discussion} interprets the findings, and Section~\ref{sec:conclusion} closes with the writer-identification extension.
